\documentclass[acmsmall,screen,review,anonymous]{acmart}

\setcopyright{none}
\settopmatter{printacmref=true}
\renewcommand\footnotetextcopyrightpermission[1]{}

\acmConference[miniKanren 2026]{The 2026 miniKanren and Relational Programming Workshop}{August 24, 2026}{Indianapolis, Indiana, USA}
\acmYear{2026}
\copyrightyear{2026}

\title{Proflog: A Tableau-Based Logic Programming Kernel in miniKanren}

\begin{document}

\begin{abstract}
We report on an implementation of Melvin Fitting's Proflog, a logic programming
language based on the tableau proof procedure.  The system extends a first-order
tableau prover with equality, written in Clojure's \texttt{core.logic}, with
Fitting's Procedure Call rule.  Formulae thereby acquire a programmatic
character: evaluating a program is a matter of closing the relevant
tableau(s).  Following alphaLeanTAP, successful evaluation also produces an
explicit proof trace.  We describe the proof kernel, the surface-to-kernel
translation, and worked executions of Fitting's reference programs.
\end{abstract}

\ccsdesc[500]{Theory of computation~Logic and verification}
\ccsdesc[300]{Theory of computation~Automated reasoning}
\ccsdesc[300]{Software and its engineering~Constraint and logic languages}

\keywords{miniKanren, logic programming, tableaux, Proflog, proof search}

\maketitle

\section{Introduction}

Fitting proposed Proflog as an exploration in alternative bases for logic
programming~\cite{fitting1994tableaux}.  By comparison, Prolog
operationalizes Horn clauses by resolution; Proflog instead manipulates
formulae of first-order logic with equality by a tableau proof procedure,
augmented with a mechanism for recursive procedure calls.  The core tableau
procedure employs standard branch closure, alpha/beta decomposition of
connectives, and gamma/delta treatment of quantifiers.  Equality is defined
over a weak Herbrand setting, in which equal terms must be syntactically
identical after branch substitutions are applied.  The source code of a Proflog
program is a series of clauses, each an implication between a formula \(\phi\)
and an atomic formula \(R\):
\[
  R(x_1,\ldots,x_n) \leftarrow \phi(x_1,\ldots,x_n).
\]
``Programmaticity'' is introduced via the Procedure Call rule, which allows
clauses to compose.  During an attempt to close a tableau, a saved literal
whose relation is defined by the program may open a subsidiary tableau for the
corresponding clause body; the subsidiary proof then contributes to closure of
the original branch.  Figure~\ref{fig:layers} summarizes the implementation
layers around this kernel rule.

\begin{figure}[t]
  \centering
  \begin{tabular}{c}
    \fbox{\begin{minipage}{0.78\linewidth}\centering
      \texttt{proflog} REPL and paper examples
    \end{minipage}} \\
    \(\downarrow\) \\
    \fbox{\begin{minipage}{0.78\linewidth}\centering
      Prefix frontend: \texttt{language}, \texttt{proflog}, \texttt{q},
      \texttt{run}
    \end{minipage}} \\
    \(\downarrow\) \\
    \fbox{\begin{minipage}{0.78\linewidth}\centering
      Compiled language and clause table
    \end{minipage}} \\
    \(\downarrow\) \\
    \fbox{\begin{minipage}{0.78\linewidth}\centering
      Query layer: success, failure, answer export, status
    \end{minipage}} \\
    \(\downarrow\) \\
    \fbox{\begin{minipage}{0.78\linewidth}\centering
      Proof-profile dispatcher
    \end{minipage}} \\
    \(\downarrow\) \\
    \fbox{\begin{minipage}{0.78\linewidth}\centering
      Tableau kernel with equality and Procedure Call
    \end{minipage}}
  \end{tabular}
  \caption{Reviewer-visible Proflog layers from surface programs to the
  proof-producing kernel.}
  \Description{A vertical stack showing the Proflog REPL and examples,
  prefix frontend, compiled language and clause table, query layer,
  proof-profile dispatcher, and tableau kernel.}
  \label{fig:layers}
\end{figure}

This changes the usual programming/proving boundary.  In Proflog, a program is
not a collection of Horn clauses with procedural negation as failure.  It is a
finite collection of first-order definitions, at most one per relation symbol,
whose queries are interpreted by two proof searches: a query succeeds when the
negated query closes, fails when the query itself closes, and is unresolved when
neither direction closes within the admitted search.  This mirrors Fitting's
supervaluation semantics and makes some operational warnings visible, especially
the warning that factoring a formula into an auxiliary relation may change the
program.

The implementation reported here is written in Clojure and \texttt{core.logic}
in the miniKanren tradition~\cite{byrd2009relational,corelogic}.  Its immediate
prover lineage is alphaLeanTAP, a relational tableau prover for first-order
classical logic~\cite{near2008alphaleantap}, and a Clojure/core.logic.nominal
implementation reference by Amin~\cite{aminLeantap}.  Proflog extends that
lineage in the direction Fitting sketched: first-order proof search becomes a
program evaluator, and program evaluation returns proof traces whose step tags
record the tableau work performed.

\section{Implementation}

\subsection{Proof Kernel}

The kernel is a relation over proof states.  Its central entry point is
\texttt{prove-stateo}; public calls use \texttt{kernel/prove} for direct
formulae and \texttt{kernel/prove-program} for program-bearing formulae.  A
branch state contains the focused formula, pending formulae, saved literals,
the lexical environment for bound variables, gamma-introduced proof variables,
delta parameters, an explicit equality substitution, delayed disequalities, the
compiled program, fuel, and the proof term being constructed.

The source does not package this state in a record; the slots are the argument
vector of the kernel relation itself:

\begin{verbatim}
(defn prove-stateo
  [fml unexpanded lits env proof-vars
   sigma sigma-out neqs neqs-out
   prog gamma-terms fuel proof]
  ...)
\end{verbatim}

The equality design is deliberately explicit.  Positive equality extends a
free-constructor substitution; negative equality is retained in a symbolic
disequality store until later bindings force closure.  This keeps equality,
disequality, and procedure-call interaction visible in proof objects.  A small
closure proof, for example, has the shape:

\begin{verbatim}
(conj (savefml (close)))
\end{verbatim}

An equality/disequality proof exposes the order of work:

\begin{verbatim}
(conj (neq-store (eq-step (eq-bind) (neq-close))))
\end{verbatim}

Procedure calls add \texttt{pos-call} and \texttt{neg-call} proof steps.  A
positive call closes the relation body.  A negative call closes the normalized
negation of that body.  Both calls run under the same kernel relation rather
than through host-language evaluation of the source program.

\subsection{Surface Syntax and Descent}

The implemented frontend is a Clojure-readable prefix DSL that preserves
Fitting's clause organization.  The intended handwritten notation is:

\begin{verbatim}
p(x) :- x = a.

only-zero(x) := forall y. (x != y or y = zero).
zero-only(x) :- only-zero(x).
\end{verbatim}

The frontend form uses \texttt{|-} for a real Proflog clause and \texttt{:=}
for a frontend-only abbreviation:

\begin{verbatim}
(def peano-language
  (pf/language
    (constants zero)
    (functions (s 1))
    (relations (zero-only 1))))

(def zero-only-program
  (pf/proflog peano-language
    (:= (only-zero x)
      (forall [y]
        (or (!= x y)
            (= y zero))))
    (|- (zero-only x)
        (only-zero x))))
\end{verbatim}

After macro expansion and compilation, the kernel sees tagged formula data.  The
simple query \texttt{p(a)} descends schematically to:

\begin{verbatim}
(pos (app p (app a)))
\end{verbatim}

The inlined \texttt{zero-only/1} body descends schematically to:

\begin{verbatim}
(forall
  (tie y
    (or
      (neq (var x) (var y))
      (eq (var y) (app zero)))))
\end{verbatim}

The public query API then probes the kernel in two directions:

\begin{verbatim}
(query/query-succeeds program query n fuel)
(query/query-fails    program query n fuel)
\end{verbatim}

\texttt{query-succeeds} asks the kernel to close the negated query;
\texttt{query-fails} asks it to close the query itself.  Open-answer examples
use \texttt{pf/run}, which binds visible answer variables in the style of
miniKanren's \texttt{run}, translates the query to the same backend formula,
and exports bindings from the proof-backed answer layer.

\subsection{Fitting's Reference Programs}

Fitting's first reference program defines even and odd natural numbers over the
language with constant \texttt{zero} and unary successor \texttt{s}:

\begin{verbatim}
even(x) :- x = zero
        or exists y. (x = s(y) and odd(y)).

odd(x) :- forall y. (even(y) -> x != y).
\end{verbatim}

The implemented frontend form is:

\begin{verbatim}
(pf/proflog p1-language
  (|- (even x)
    (or (= x zero)
        (exists [y]
          (and (= x (s y))
               (odd y)))))
  (|- (odd x)
    (forall [y]
      (implies (even y)
               (!= x y)))))
\end{verbatim}

The second reference program is a one-or-two-token Nim game.  A position wins
when there is a legal move to a position that does not win:

\begin{verbatim}
win(x) :- exists y.
            ((x = s(y) or x = s(s(y))) and not win(y)).
\end{verbatim}

The implementation keeps the move test inlined, as in Fitting's presentation:

\begin{verbatim}
(pf/proflog p2-language
  (|- (win x)
    (exists [y]
      (and (or (= x (s y))
               (= x (s (s y))))
           (not (win y))))))
\end{verbatim}

Table~\ref{tab:fitting-results} summarizes the focused Fitting suite results.
It covers forward success, forward refutation, answer synthesis, partial
synthesis, and unresolved classification.

\begin{table}[h]
  \caption{Representative Fitting-program outcomes}
  \label{tab:fitting-results}
  \begin{tabular}{lll}
    \toprule
    Case & Mode & Outcome \\
    \midrule
    \texttt{even(0)} & forward success & succeeds \\
    \texttt{odd(s(0))} & forward success & succeeds \\
    \texttt{odd(0)} & forward refutation & fails \\
    \texttt{win(4)} & forward success & succeeds \\
    \texttt{win(3)} & forward refutation & fails \\
    \texttt{move(1,0)} & auxiliary forward success & succeeds \\
    \texttt{move(0,1)} & auxiliary forward refutation & fails \\
    factored \texttt{win(1)} & status classification & unresolved \\
    \texttt{append(x,y,[a,b,c])} & answer synthesis & target found \\
    \texttt{reverse([a,b,c,a],cons(a,r))} & partial synthesis & target found \\
    \bottomrule
  \end{tabular}
\end{table}

A representative full output row for Fitting's first program is:

\begin{verbatim}
[{:description
  "P1 proves odd(s(0)) through the forall-based odd clause",
  :family :p1,
  :outcome :succeeds,
  :classification nil,
  :max-fuel nil,
  :id :p1-odd-1-succeeds,
  :kind :query,
  :proof-count 1,
  :expected :succeeds,
  :proof-steps
  [neg-call
   witness
   conj
   eq-step
   par-bind
   pos-call
   split
   free-close
   witness
   conj
   eq-step
   decompose
   args
   par-bind
   pos-call
   univ
   split
   neg-call-guarded-alt
   guarded-alt
   guarded-neg-alt-saturated
   guarded-scope-done
   guard-eq
   eq-bind
   guard-saturation-done
   guarded-call-seq-done
   guarded-seq-done
   refl-close],
  :proof-root neg-call}]
\end{verbatim}

Most tags in this trace are ordinary tableau or equality work, such as
\texttt{witness}, \texttt{conj}, \texttt{split}, \texttt{univ},
\texttt{eq-step}, and \texttt{refl-close}.  The tags that are specific to
Proflog relative to alphaLeanTAP are the procedure-call and guarded-clause
steps: \texttt{neg-call}, \texttt{pos-call},
\texttt{neg-call-guarded-alt}, \texttt{guarded-alt}, and the
\texttt{guarded-*} saturation and sequence tags.  These are the proof-trace
footprints of Fitting's Procedure Call rule.

The factored Nim case is included because it demonstrates that the evaluator is
not merely computing the Nim recurrence in the host language.  If the program
introduces an auxiliary relation

\begin{verbatim}
win(x) :- exists y. (move(x, y) and not win(y)).
move(x, y) :- x = s(y) or x = s(s(y)).
\end{verbatim}

then \texttt{move/2} itself is decidable on ground calls, but \texttt{win(1)}
remains unresolved in the tested kernel setting.  This is the operational
boundary Fitting warned about: a logically tempting refactoring changes how
procedure calls are admitted and discharged.  In the inlined program, the two
move alternatives remain ordinary equality and disjunction work inside the
subsidiary tableau opened for \texttt{win/1}.  In the factored program, the
same alternatives sit behind a second procedure boundary.  Proflog's Procedure
Call rule is not a general unfold/fold equivalence principle; it admits calls
when a saved literal and the current branch state make the corresponding clause
body usable.  Moving formula structure into a new relation can therefore alter
which subsidiary tableaux are available at the moment a branch needs to close.
The example is small, but it exposes the larger point: Proflog programs are not
only extensional formula collections, but operational tableau objects.

\subsection{Proof Profiles}

Fitting emphasizes that Proflog gains expressiveness by using tableaux as the
programming substrate, but that this substrate also inherits the search costs
of first-order proof.  The implementation therefore includes proof profiles:
small, proof-producing dispatch layers for fragments where a specialized
kernel can close a branch without changing the meaning of a successful proof.
The current profiles recognize propositional formulae, equality-free
first-order formulae, and equality-dominated finite fragments used by several
paper examples.

Profiles are not a separate evaluator for Proflog source programs.  They are
entered only after the surrounding branch has been isolated from active
procedure calls and the relevant equality state, and their subproof is recorded
under an explicit profile tag in the proof trace.  This keeps the main
Procedure Call semantics intact while allowing the implementation to reuse
Fitting's recognition that different fragments of the logic have different
operational envelopes.  The same architecture can be reused to extend and
experiment with future Proflog implementations: a profile can be added for a
well-understood fragment, tested as a proof-producing branch rule, and removed
without changing the surface language.

\section{Discussion}

Proflog is of interest because it gives another concrete point in the space of
connections between logic and computation.  The current implementation
demonstrates that the core operational semantics of the tableau proof
procedure, augmented with the Procedure Call rule, can be defined as a pure
relation in the miniKanren style.  Proof search, program evaluation, and
proof-trace construction are furthermore extensible for future experimentation
via the profile architecture.

The implementation is faithful to Fitting's specification in both its
strengths and its limitations.  Fitting's semantics require query failure to be
proved, rather than assumed by search exhaustion.  Undefined and unresolved
cases are first-class outcomes.  Refactorings that would be harmless in a
purely extensional reading can change operational behavior when they introduce
auxiliary procedure calls.

Performance remains an area for future work: recursive examples remain
search-control frontiers, and several successful answer and partial-synthesis
tests are intentionally kept out of the default fast path.  However, a gap in
the literature has been covered: Fitting's theoretical Proflog can now be
considered mechanized as a proof-producing miniKanren-style kernel, and its
non-trivial examples can be evaluated through that kernel with correctness
results visible at the proof boundary.

\begin{acks}
Acknowledgments are omitted for anonymous review.
\end{acks}


\bibliographystyle{ACM-Reference-Format}
\bibliography{references}

\end{document}
